% =====================================================================
% Bagian 1 — Pendahuluan  (target 1,25 halaman)
% MASALAH dan PRODUK wajib konsisten dengan concept paper penyisihan.
% Urutan kontribusi SENGAJA diubah: yang paling kuat buktinya di depan.
% =====================================================================

\section{Pendahuluan}

Kebakaran hutan dan lahan (karhutla) gambut di Sumatra dan Kalimantan bukan
sekadar bencana musiman melainkan persoalan struktural yang berdampak luas pada
kesehatan masyarakat, ekonomi, dan hubungan diplomatik regional. Kejadian ini
kronis sejak setidaknya 1997: 2015 tercatat sebagai salah satu yang terburuk
dengan kerugian ekonomi US\$28 miliar dan enam kejadian 2004--2015 mencapai
kumulatif US\$93,9 miliar \cite{worldbank2019}; 2019 menghanguskan sekitar
942.000 hektare dengan kerugian sedikitnya US\$5,2 miliar dan lebih dari
900.000 laporan gangguan pernapasan, serta menutup 24.773 sekolah dan
menghentikan pembelajaran 4.692.537 siswa \cite{unicef2019}. Pola berulang pada
1997, 2015, dan 2019 menunjukkan akar persoalannya adalah \textbf{keterlambatan
deteksi, bukan ketiadaan kemauan penanganan}: ketiganya baru mendapat perhatian
luas setelah asap menyelimuti wilayah lintas negara. Karena biaya pemadaman
meningkat tajam seiring meluasnya area yang terbakar sebelum terdeteksi,
penajaman deteksi pada fase paling awal adalah titik ungkit dengan dampak
terbesar.

\subsection{Kesenjangan operasional dan solusi yang diusulkan}

Infrastruktur pemantauan yang tersedia luas berbasis \emph{hotspot} satelit
seperti FIRMS dan SiPongi. Sistem ini efektif sebagai pemicu awal berbiaya
rendah pada skala luas, tetapi resolusi spasial dan temporalnya tidak memadai
untuk mengonfirmasi titik api kecil yang tersamar asap tebal atau kanopi gambut.
Muncullah kesenjangan operasional antara pemicu kasar berskala satelit dan
kebutuhan verifikasi visual beresolusi tinggi sebagai dasar pengerahan tim.
Wawancara dengan narasumber ahli (Bagian~5) menegaskan bentuk kesenjangan ini:
drone memang telah banyak dipakai, tetapi pemanfaatannya untuk karhutla terbatas
pada pengambilan citra berkala \emph{setelah} kebakaran diketahui terjadi,
dengan penyebaran api kemudian diprediksi manual. Fase yang belum tergarap
justru fase sebelum konfirmasi.

Solusi yang kami tawarkan adalah \textbf{sistem deteksi dini karhutla gambut
berjenjang} yang mengisi kesenjangan tersebut melalui konfirmasi visual
beresolusi tinggi berbasis drone. Inti kontribusinya adalah model fusi
RGB--termal berbasis \emph{cross-attention} dengan \emph{modality dropout} agar
tetap andal ketika salah satu sensor tidak tersedia atau terganggu, kondisi
lazim pada asap pekat maupun operasi malam. Keluarannya adalah \emph{alert} dini
dan peta risiko yang dapat langsung ditindaklanjuti, disajikan melalui konsol
operator yang berjalan sebagai sistem terpasang, bukan mockup.

Sasarannya dua segmen. Pertama, instansi penanggulangan bencana dan kehutanan
(Manggala Agni, BNPB, dinas kehutanan daerah), yang memperoleh dasar objektif
untuk memprioritaskan pengerahan tim. Kedua, pemegang konsesi perkebunan sawit
dan hutan tanaman industri di lanskap gambut, yang memikul tanggung jawab mutlak
atas kebakaran di arealnya tanpa perlu pembuktian kesalahan menurut Pasal~88
UU~32/2009 \cite{uu32_2009}, sehingga adopsi menjawab insentif kepatuhan dan
bukan sekadar itikad baik. Rincian titik nyeri kedua segmen, beserta sumbernya,
ada pada Lampiran~\ref{app:pengguna}.

\subsection{Kontribusi}

Kontribusi kami urutkan menurut kekuatan bukti yang benar-benar diperoleh,
bukan menurut ambisi awal.

\textbf{(1)~Protokol evaluasi ketahanan berbasis kurva degradasi kontinu.}
Ketahanan diukur sebagai kurva atas enam level degradasi tersintesis dari citra
nyata, bukan uji biner. Protokol ini menghasilkan dua temuan yang tidak akan
muncul dari uji biner: evaluasi pada kondisi bersih terbukti \emph{tidak memiliki
daya pisah} pada tolok ukur ini, ketiga konfigurasi modalitas mencapai batas
atas yang sama persis pada partisi uji, dan peringkat ketahanan antar-modalitas
\emph{berbalik} antar-partisi.
\textbf{(2)~Perluasan \emph{modality dropout} menjadi \emph{reliability gating}},
dari regularisasi pasif menjadi \emph{gating} adaptif yang disupervisi-diri oleh
label degradasi hasil pipeline augmentasi sehingga tidak memerlukan anotasi
tambahan, dengan skor keandalan per modalitas yang sampai ke operator sebagai
informasi keputusan. Kami melaporkan secara terbuka bahwa \emph{gating}
\textbf{tidak meningkatkan akurasi klasifikasi}.
\textbf{(3)~Ablation lokalisasi supervisi-lemah versus supervisi-penuh}, salah
satunya tanpa label piksel sama sekali, syarat praktis untuk ekspansi ke lanskap
yang belum teranotasi.
\textbf{(4)~Sistem terpasang ujung-ke-ujung} dengan integrasi hotspot satelit
langsung, disertai audit kejujuran representasi: komponen simulasi berlabel, dan
besaran yang belum terukur ditampilkan kosong alih-alih diisi angka rekaan.

\subsection{Batasan penelitian}
\label{sec:batasan}

Agar seluruh klaim dapat diverifikasi dalam jendela waktu kompetisi, penelitian
ini dibatasi dan pembatasannya kami nyatakan di muka. \textbf{(i)}~Tidak ada data
gambut tropis: pelatihan dan evaluasi memakai FLAME~2 beserta anotasi piksel
RFFNet (kebakaran bernyala, hutan pinus Arizona), sehingga sistem yang dirancang
untuk rezim membara ini divalidasi pada proksi terukur dan belum tervalidasi
untuk gambut. \textbf{(ii)}~Kondisi lingkungan pada proses prediksi dibatasi pada kondisi yang
direpresentasikan dalam dataset FLAME~2, yaitu lingkungan kebakaran hutan yang
diamati melalui citra udara menggunakan UAV. Prediksi dilakukan berdasarkan
karakteristik visual yang terdapat pada citra, sehingga kondisi lingkungan
seperti cuaca, suhu, kelembapan, kecepatan angin, dan karakteristik lingkungan
lainnya tidak digunakan sebagai variabel masukan secara eksplisit. Model
ditujukan untuk melakukan deteksi berdasarkan informasi visual pada citra dan
tidak dimaksudkan untuk memprediksi kejadian kebakaran berdasarkan faktor
lingkungan secara langsung. \textbf{(iii)}~Tidak ada pengujian pada perangkat
tepi fisik; klaim latensi terbatas pada proksi CPU x86. \textbf{(iv)}~Tidak ada
baseline eksternal; seluruh perbandingan bersifat internal antar-konfigurasi
modalitas dari model yang sama. \textbf{(v)}~Satu benih acak tanpa pengulangan,
sehingga selisih kecil berada dalam batas derau; setiap angka utama disertai
selang kepercayaan dan uji signifikansi berpasangan. \textbf{(vi)}~Lingkup
dibatasi pada Lapisan~0--1.
